\section{Evaluation}
\label{sec:evaluation}

Across V6, V7, and V8, all 45 planned positive rows are admitted. No measured
row times out, reaches critical pressure, or grows swap. Each protocol produces
one response SHA-256 across all three policies, and every candidate row observes
nonzero empty-slot recovery. Paired comparisons use the deterministic 95\%
bootstrap confidence interval defined in \cref{sec:methodology}. We first report the negative precursor, then answer
the five research questions.

\subsection{Negative precursor: hit rate is not the objective}

\Cref{fig:gate1} shows the fixed-cache sweep. Hit rate rises monotonically from
48.98\% at 16 slots to 90.13\% at 96 slots. Decode throughput moves in the
opposite direction, while minimum free memory falls from 51\% to 32\% and peak
compressor growth rises from 714.8\mib to 5,830.0\mib. The matched 16-to-96
slowdowns are 10.72\%, 10.84\%, 10.54\%, 9.45\%, and 8.94\%. Their paired
median is 10.54\%, with a bootstrap interval of 8.94--10.84\%.

The tradeoff is stable but below the preregistered 15\% controller gate. Gate 1
is therefore rejected. This matters to the positive result: \system is not a
post-hoc cache size selected from the sweep. It changes memory ownership and
validity semantics, and its tests use new prompts and experiment identifiers.

\subsection{RQ1--RQ3: Qwen throughput and physical pressure}

\begin{table}[ht]
\centering
\caption{Qwen3.6 policy results. Values are medians of five admitted measured
runs per arm. Footprint is peak \code{ri\_phys\_footprint}; compressor is the
per-run peak delta above the pre-run snapshot.}
\label{tab:qwen}
\small
\begin{tabular}{llrrrrr}
\toprule
Protocol & Policy & \multicolumn{1}{c}{tok/s} & \multicolumn{1}{c}{Hit} &
\multicolumn{1}{c}{Footprint} & \multicolumn{1}{c}{Compressor} &
\multicolumn{1}{c}{Min free} \\
\midrule
\multirow{3}{*}{4\,GiB co-tenant}
 & Upstream-16 & 13.025 & 47.91\% & 1.374\,GiB & 0.675\,GiB & 57\% \\
 & Upstream-96 & 12.116 & 90.81\% & 6.705\,GiB & 5.890\,GiB & 38\% \\
 & \system-96/16 & \textbf{14.136} & 86.83\% & 4.584\,GiB & 0.840\,GiB & 44\% \\
\midrule
\multirow{3}{*}{No pressure worker}
 & Upstream-16 & 14.520 & 49.82\% & 1.386\,GiB & 0.000\,GiB & 57\% \\
 & Upstream-96 & 14.293 & 90.58\% & 6.721\,GiB & 5.991\,GiB & 38\% \\
 & \system-96/16 & \textbf{16.807} & 87.50\% & 4.594\,GiB & 0.700\,GiB & 44\% \\
\bottomrule
\end{tabular}
\end{table}

Under the 4\,GiB co-tenant, \system gains 7.25\% paired-median throughput over
upstream-16 and 13.67\% over upstream-96. The interval against upstream-16
spans -5.02\% to +24.08\%; against upstream-96 it is +5.33\% to +33.91\%.
The candidate wins nine of ten pressure-backed comparisons and its only loss is
5.02\%, inside the frozen fallback gate. It retains 38.92 percentage points
more hit rate than upstream-16.

Against upstream-96, peak physical footprint falls from 6.705 to 4.584\,GiB,
a 31.64\% reduction. Peak compressor growth falls from 5.890 to 0.840\,GiB, a
5.05\,GiB benefit. \system is not the smallest-memory policy: upstream-16 has
a 1.374\,GiB footprint. The candidate dominates upstream-96 in the measured
throughput-footprint plane. Against upstream-16 it trades a larger footprint
for higher median throughput and 38.92 percentage points more locality.

Without a synthetic pressure worker, the same policy gains 16.91\% over
upstream-16 (interval 4.56--38.79\%) and 17.59\% over upstream-96 (interval
4.79--29.15\%). It wins all ten matched comparisons. Fixed-96 footprint falls
31.65\%, and compressor growth falls by 5.29\,GiB. The effect therefore does
not require the controlled co-tenant, although the absolute rates differ
because the host begins from a different memory state.

\begin{figure}[tb]
  \centering
  \includegraphics[width=0.98\linewidth]{figures/elasticuma-main-results.pdf}
  \caption{Qwen3.6 policy comparison. With the co-tenant, median throughput is
  14.136\tps for \system versus 13.025 and 12.116\tps for upstream-16 and
  upstream-96. Without it, the values are 16.807, 14.520, and 14.293\tps.
  Fixed-96 footprint falls by about one third under both protocols. Bars are arm
  medians; points are all measured runs.}
  \label{fig:qwen-main}
\end{figure}

\begin{figure}[tb]
  \centering
  \includegraphics[width=0.90\linewidth]{figures/elasticuma-paired-gains.pdf}
  \caption{Matched Qwen decode gains against upstream-16 and upstream-96.
  Paired medians are 7.25\% and 13.67\% with the co-tenant and 16.91\% and
  17.59\% without it. Horizontal intervals are deterministic bootstrap 95\%
  intervals from five matched repetitions; the shaded region marks slowdowns.}
  \label{fig:paired}
\end{figure}

\subsection{RQ4: recovery is exercised, not inferred}

Every Qwen candidate row records volatile transitions, relocks, and a nonzero
empty-recovery count. Cumulative recovered bytes range from 9.090 to
10.758\,GiB per pressure-backed run and from 7.907 to 10.331\,GiB per
no-pressure run. These totals exceed instantaneous footprint because a cold
slot can be reclaimed and reloaded repeatedly.

All policies within V6 produce one identical response hash; V7 produces a
second within-protocol hash. These matching hashes demonstrated response-text
equality under greedy decoding despite real empty recovery. The current upstream CLI does not
expose generated token IDs, so we do not claim token-ID equality. The evidence
is stronger than a no-reclamation test but narrower than a full numerical trace
comparison.

\FloatBarrier
\pagebreak[4]

\subsection{RQ5: transfer to Gemma 4}

\begin{table}[ht]
\centering
\caption{Gemma 4 second-architecture validation without a pressure worker.
Values are medians of five admitted measured runs per arm.}
\label{tab:gemma}
\small
\begin{tabular}{lrrrrr}
\toprule
Policy & \multicolumn{1}{c}{tok/s} & \multicolumn{1}{c}{Hit} &
\multicolumn{1}{c}{Footprint} & \multicolumn{1}{c}{Compressor} &
\multicolumn{1}{c}{Min free} \\
\midrule
Upstream-16 & 16.394 & 65.52\% & 1.986\,GiB & 0.000\,GiB & 69\% \\
Upstream-96 & 18.742 & 98.40\% & 8.538\,GiB & 5.114\,GiB & 50\% \\
\system-96/16 & \textbf{23.820} & 95.79\% & 5.527\,GiB & 1.073\,GiB & 55\% \\
\bottomrule
\end{tabular}
\end{table}

\FloatBarrier

Gemma changes layer count, expert count, tokenizer, routing, and per-slot size.
To test architecture transfer without tuning, we preserve the Qwen 96/16
residency policy. Against the pinned upstream-16 and upstream-96 policies
\cite{patel2026slipstream}, \system wins every matched comparison. Paired-median
gains versus upstream-16 and upstream-96 are 45.30\% (interval 23.45--51.44\%)
and 26.01\% (interval 11.88--28.55\%), respectively. Relative to fixed-96, peak
footprint falls 35.26\%, from 8.538 to 5.527\,GiB; compressor growth falls from
5.114 to 1.073\,GiB. Hit rate remains 30.28 percentage points above the small
cache. All candidate rows exercise exact empty recovery (2.421--2.868\,GiB of
cumulative reload traffic), and no row grows swap.

\FloatBarrier

\begin{figure}[ht]
  \centering
  \includegraphics[width=0.98\linewidth]{figures/elasticuma-cross-model.pdf}
  \caption{No-pressure results across two MoE architectures. Absolute token
  rates are not a cross-model benchmark because prompts and generation lengths
  differ. Within each model, \system has the highest median decode rate and a
  materially lower footprint than upstream-96.}
  \label{fig:cross-model}
\end{figure}

\FloatBarrier

\subsection{Rejected alternatives and mechanism ablation}

\Cref{tab:ablations} records the measured failure criterion for each rejected
mechanism. This is a mechanism-screening record, not a quantitative aggregate
across incompatible prototypes.

\begin{table}[ht]
\centering
\caption{Alternatives rejected before the final per-slot, layerwise design.}
\label{tab:ablations}
\begin{tabularx}{0.98\linewidth}{>{\bfseries}l X}
\toprule
Alternative & Observed failure \\
\midrule
\code{F\_NOCACHE} & Advice changed file-cache treatment but did not release the anonymous slot allocation. \\
Direct mmap & Reduced measured model RSS from 3.38 to 1.92\,GiB in a probe, but was approximately 24 times slower than the cached positional-read path. \\
Raw/compressed MTLIO & Direct-storage variants missed frozen latency gates; lossless CPU decode reached only 0.46--2.85\,GiB/s. \\
One arena per layer & Coarse purgeability collapsed Qwen effective hit rate to 1.3\% and roughly halved throughput. \\
Post-prefill-only policy & Correctly marked cold slots but ran after the full cache had already materialized, failing the footprint gate. \\
\bottomrule
\end{tabularx}
\end{table}

\FloatBarrier

Only Metal-owned per-expert buffers, selective hot/cold state, and transitions
inside the synchronized layer loop satisfy both validity and footprint gates.

\subsection{Dense Qwen3.8 diagnostic}

We also run one canonical Qwen3.8-27B Q4 dense checkpoint with MLX
\cite{apple2026mlx}. One warmup and one 128-token measured row complete at
11.36\tps. The sampled peak physical footprint is 17.96\,GiB; MLX reports
18.60\,GiB peak memory. Minimum free memory reaches 20\%, and swap grows by
151\mib. A later row reaches
19\% free memory and is terminated by the frozen guard. Because the planned
five measured rows do not complete, this diagnostic is not admitted. We do not
use it as a speed, quality, or dense-versus-MoE comparison.

\FloatBarrier
